\section{Approach C: cluster, then walk backwards from the ideal key}
\label{sec:keys}
Approach C changes the lookup structure. Instead of repeatedly intersecting document masks,
it asks: \emph{which stored documents have this particular key?} Start from the signs the query
prefers, then consider other keys in increasing mismatch penalty.

The primary method here still starts with the same selected clusters. A global hash table is
the special case $C=P=1$. We do not silently switch between clustered and global formulas.

\fig{keys}{Approach C. Each cluster has a key directory and postings. Empty lookups still consume work.}{1}

\subsection{Full keys and short keys are different choices}
Let $I_h$ be the fixed set of $h$ lookup coordinates, often the first $h$ positions. A
\textbf{full key} has $h=d$. Its weighted penalty determines the complete binary score.
A \textbf{short key} has $h<d$. It ranks a prefix, then the full $d$-bit score ranks the
retrieved candidates.

A full-key search can be exact if it enumerates keys in true query-weighted order and finishes
the result-boundary tie. A short-key candidate quota does not have that guarantee. It can still
work well, and can also be made exact with a valid full-score stopping bound. These are
properties of the stopping rule, not of the name ``walk backwards.''

\subsection{The data structure behind a key lookup}
For each cluster, a directory maps a key to a posting range containing the local row references
of documents with that key. There may be many documents per key. Finding a directory entry
without finding its document rows is not enough to produce results.

\begin{lstlisting}[caption={Build a clustered key index.}]
BuildKeyIndex(cluster_codes, lookup_coordinates):
    for each cluster j:
        groups = group local row references by their lookup key
        concatenate the groups into a posting array
        for each occupied key:
            directory[j][key] = (posting_start, posting_count)
    return directories, postings
\end{lstlisting}

A \textbf{dense directory} has an offset for every possible key. A \textbf{sparse hash table}
only stores occupied keys but adds key, slot and load-factor overhead. A sorted-key array is
another option: it can use less overhead, but a lookup uses comparisons or another search
structure instead of an assumed constant-time hash access.

\subsection{Do not index fixed routing bits as if they were still free}
If sign-prefix routing fixes some lookup coordinates inside cluster $j$, those coordinates
cannot vary within it. Let $r_j$ be the number of lookup coordinates fixed by that cluster,
and let
\begin{equation}
h_j=h-r_j,\qquad U_j\leq\min(n_j,2^{h_j}).
\label{eq:free-key}
\end{equation}
$U_j$ is the number of occupied free keys. For first-$r$-bit routing and a first-$h$-bit lookup
with $h\geq r$, $h_j=h-r$. The fixed coordinates add a \emph{constant} penalty for that cluster;
they do not create $2^r$ additional keys inside it. Ordinary k-means clusters do not generally
fix any sign coordinate, so this simplification does not follow from their cluster ID.

\begin{examplebox}{Two fixed bits and a two-bit local key}
Suppose a four-bit code is routed by its first two bits. Inside cluster \texttt{10}, codes
\texttt{1010}, \texttt{1011}, \texttt{1000}, \texttt{1001} need only the local keys
\texttt{10}, \texttt{11}, \texttt{00}, \texttt{01}. A lookup for the full prefix must include
the cluster's constant routing penalty when it is compared with a different cluster's stream.
The number of local possibilities is four, not sixteen.
\end{examplebox}

\subsection{Index storage: name every array}
Use the common baseline row order and ID mapping. Let $b_{\rm ref}$ be the bytes for a local
posting reference, $b_k(h_j)$ the bytes storing a key, and $b_{\rm slot}$ the remaining per-slot
metadata. With a chosen hash load factor $0<\alpha<1$, a reference sparse-key layout uses
\begin{equation}
M_{\rm keydir}=\sum_j\ceil{U_j/\alpha}
                 \left[b_k(h_j)+2b_o+b_{\rm slot}\right],
\label{eq:hash-storage}
\end{equation}
where the two offsets represent start and length. Field alignment or different slot layouts
can be substituted explicitly. The additional posting references cost $N b_{\rm ref}$.
Therefore
\begin{equation}
\boxed{M_{C,\rm index}=\Mbase+N b_{\rm ref}+M_{\rm keydir}.}
\label{eq:keys-storage}
\end{equation}
This chosen layout deliberately counts both the common external-ID mapping and the extra
row-reference permutation. If the build reorders the actual codes and IDs by key within each
cluster, postings can become direct contiguous row ranges. Then the extra $N b_{\rm ref}$ can
be removed. That is an explicit layout change, not permission to omit all document IDs.

For a dense offset directory instead of a hash table,
\[
M_{\rm dense-dir}=\sum_j(2^{h_j}+1)b_o.
\]
Exponential key space makes dense directories expensive. Sparse directories avoid storing
empty keys, but do not make probing empty keys free.

\subsection{How to enumerate all flip combinations in weighted order}
Let the free-coordinate weights be sorted in ascending order,
$a_{(1)}\leq\cdots\leq a_{(h_j)}$. Flipping a subset $S$ costs
$\sum_{i\in S}a_{(i)}$. Merely flipping the first bit, then the first two, then the first three
misses other subsets. We need a stream containing \emph{all} subsets in order.

The following priority-queue construction has at most two children per nonempty subset.
It starts with the empty subset, then the subset containing coordinate 1.
\begin{lstlisting}[caption={A complete sorted-subset stream, not one chain of flips.}]
SortedFlipSets(weights):
    sort weights ascending, remembering original coordinate IDs
    emit (empty subset, penalty 0)
    if there are no free coordinates: stop
    queue.push(subset={1}, last=1, cost=weights[1])
    while queue is not empty:
        (S, i, cost) = queue.pop_smallest_cost()
        emit (S, cost)
        if i < number_of_weights:
            queue.push(S union {i+1}, i+1, cost + weights[i+1])
            queue.push((S minus {i}) union {i+1}, i+1,
                       cost - weights[i] + weights[i+1])
\end{lstlisting}

Both child costs are no smaller than the parent cost. Every nonempty subset has one parent:
if the coordinate before its largest coordinate is also present, remove the largest one;
otherwise replace the largest coordinate by its predecessor. This explains why the construction
visits every subset without needing a duplicate-state hash set.

A state must store enough information to construct the key. A full bitmask representation costs
$b_w\ceil{h_j/w}$ bytes per state, plus its penalty, last coordinate and queue metadata.
Copying that mask into children is work. A persistent representation can change this cost,
but it then needs its own traversal and memory accounting.

\begin{examplebox}{The full-key order can be accurate even when a greedy branch is not}
For $q=(0.6,0.5,0.4,0.3)$, the ideal code is \texttt{1111}. The first penalties include
$0$, $0.3$, $0.4$, $0.5$, $0.6$, $0.7$, and so on. The code \texttt{0111} is found by flipping
coordinate 1 at penalty $0.6$. The code \texttt{1000} has penalty $1.2$ and comes later.
Complete weighted enumeration finds the better document first, even though it disagrees on
the query's largest coordinate. It may have to check several empty keys before that hit.
\end{examplebox}

\subsection{Lookup and stopping within selected clusters}
For short keys, a simple implementation allocates candidate targets $m_j$ to probed clusters.
A more flexible implementation merges their next-prefix-key scores and draws candidates from
the best available stream. Whichever policy is used must be a parameter of the experiment;
a per-cluster target and a global target are not interchangeable.

\begin{lstlisting}[caption={Short-key candidate search with explicit per-cluster targets.}]
SearchShortKeys(q, selected_clusters, h, targets m[j], R):
    candidates = empty collection
    for each selected cluster j:
        stream = SortedFlipSets(free lookup weights in j)
        found = 0
        for (flip_set, variable_penalty) in stream:
            key = ideal_free_key(q, j) XOR flip_set
            rows = directory[j].lookup(key)     # may be empty
            keep eligible rows from this posting range
            found += number of kept rows
            append those rows to candidates
            if found >= m[j] and final prefix-penalty ties are finished:
                break
    score every candidate with the full d-bit score
    keep the best R with the same score/ID rule as A
    return best R
\end{lstlisting}

A full-key exact version instead merges the cluster streams by their \emph{total full-key}
penalty, including fixed-sign penalties. Once at least $R$ documents have been found, continue
through every key that can tie or exceed the $R$th document. This produces the exact binary
top $R$ within the selected pool. If unprobed clusters remain, global exactness additionally
requires their score bounds to be ruled out or those clusters to be searched.

A short key can also stop exactly. Let $T_h(b)$ be the prefix score and
$A_{\rm tail}=\sum_{i\notin I_h}|q_i|$. Every document in an unvisited prefix key satisfies
\begin{equation}
S(q,b)\leq T_h(b)+A_{\rm tail}.
\label{eq:prefix-bound}
\end{equation}
If the best unvisited prefix upper bound is strictly below the current worst top-$R$ score,
the result is certified. If not, continue even when a nominal candidate quota has been reached.
A larger tail makes this bound less selective. Fixed routing-coordinate contributions must be
included consistently; no coordinate may be counted twice.

\subsection{Do we need to rescore a document whose full key is known?}
Not always. With a full key and the true query weights, its complete score is already
$A-2\pi_q(\text{key})$. Every document in that key's posting range has that same binary score.
The search can offer their IDs to the result heap without unpacking and rescoring each code.
For a short key, the known prefix contribution can similarly be reused while evaluating
only the remaining coordinates.

To make that implementation choice visible, let $G_{C,j}$ be the number of score-table groups
actually evaluated per candidate in cluster $j$. The simple pseudocode above has $G_{C,j}=G$.
A full-key implementation reusing its exact key score can have $G_{C,j}=0$. A tail-only scorer
can have approximately $\ceil{d_{\rm remaining}/g}$ groups, with alignment and fixed routing
coordinates handled explicitly. It still pays for postings, key-score arithmetic, selection,
and any later float reranking. This optimization changes the cost; it is not a recall loss.

The reference storage formula retains the common packed rows. A specialized full-key layout
can also avoid a redundant per-document code copy by using its stored full keys, provided its
ID/posting and key ownership are redesigned and counted. Equation~\eqref{eq:keys-storage} is
one explicit layout, not a universal lower bound on every possible key index.

\subsection{If ``work upwards'' means removing a constraint}
There is another reasonable interpretation of the idea. Instead of flipping a bit and probing
one more exact key, remove a bit constraint. For example, \texttt{101*} means either
\texttt{1010} or \texttt{1011}. That is a set of keys. A plain exact-key hash table does not
answer this wildcard request in one ordinary lookup.

A fixed-order \textbf{prefix trie} stores paths by successive key coordinates. A node identifies
an occupied prefix. If document rows are ordered by key within a cluster, each node can store
the start and end of a contiguous posting range. Moving to its parent relaxes the last fixed
coordinate. It does not require copying every descendant's posting list into every ancestor.

\begin{examplebox}{Relaxing the key inside the example cluster}
For $\mathcal C_1=\{2,5,7,8\}$ and ideal key \texttt{1010}:
\begin{center}
\begin{tabular}{@{}lll@{}}\toprule
Constraint & Matching documents & Newly visited documents\\\midrule
\texttt{1010} & none & none\\
\texttt{101*} & 2 & 2\\
\texttt{10**} & 2, 7 & 7\\
\texttt{1***} & 2, 5, 7 & 5\\
\texttt{****} & 2, 5, 7, 8 & 8\\\bottomrule
\end{tabular}
\end{center}
The parent range contains the previous range. Scan only the newly added sibling ranges if the
implementation can retain their boundaries. Otherwise repeated scans must be counted again.
In this particular query the last coordinate is the cheapest one to relax. That is not true
for every query.
\end{examplebox}

\begin{lstlisting}[caption={A fixed-prefix relaxation variant of Approach C.}]
SearchByRelaxingPrefix(q, query_key, trie, table, R, candidate_target):
    path = follow query_key as far as the stored trie permits
    previous_range = empty
    best = empty top-R heap
    found = 0
    for node on that path, from deepest node toward the root:
        new_ranges = node.posting_range minus previous_range
        for eligible row in new_ranges:
            score its code and offer it to best
            found += 1
        previous_range = node.posting_range
        if found >= candidate_target:
            return best                       # approximate stopping rule
    return best
\end{lstlisting}
An exact version may stop early only after bounding \emph{every} unvisited region against the
current full-score heap. Simply finding enough IDs in a relaxed prefix is not such a proof.
The remaining score can still change the ordering.

Let $Z_j$ be the number of stored trie nodes and $b_{\rm trie}$ their record size, including
child links or equivalent navigation and range offsets. A simple uncompressed trie satisfies
\[
Z_j\leq 1+\sum_{u=1}^{h_j}\min(n_j,2^u).
\]
Its directory term replaces the hash-directory term by $\sum_j Z_jb_{\rm trie}$. Path compression
can reduce the actual node count. With range references, document postings are stored once;
copying a posting list into every prefix node would instead add repeated references and can
cost $O(Nh)$ reference storage.

For $u_j$ ancestor/range visits and $F_j$ total row visits, a reference cost is
\begin{equation}
T_{C,\rm trie}=T_{\rm keyprep}+
\sum_j\left[h_jc_{\rm path}+u_jc_{\rm range}+F_jc_{\rm post}
 +F_jG_{C,j}c_{\rm lut,gather}\right]+T_{\rm select/rank}.
\label{eq:trie-cost}
\end{equation}
A saved navigation path uses $O(h_j)$ records per active stream. Add the candidate heap,
gather buffers, and common query RAM. Here $F_j$ must count repeated rows if they are revisited.
The same common pipeline and final rerank terms are added for end-to-end latency.

A fixed trie can drop its last prefix coordinate cheaply. Dropping an arbitrary
query-selected low-magnitude coordinate can produce several disjoint ranges, so it does not
have the same constant-cost parent operation. A query-adaptive bit order requires additional
traversal or a different index. This is the main difference between a prebuilt fixed trie and
the query-ordered bitplane traversal.

\subsection{Count keys, postings, scores and enumeration states}
Let $H_j$ be the number of keys attempted in cluster $j$, including empty keys. Let $F_j$ be
the number of eligible posting rows retrieved and scored. Let $E_j$ be the number of
sorted-subset heap comparisons, $C_j^{\rm copy}$ the number of key-word copies, and let
$Q_j^{\max}$ be the largest enumeration-queue size. Define $F=\sum_jF_j$.

For the two-child stream, $E_j=O(H_j\log(Q_j^{\max}+1))$ comparison work, and
$C_j^{\rm copy}=O(H_j\ceil{h_j/w})$ for full-mask copies. Hashing a key also depends on its
length and table layout. An effective $c_{\rm lookup}(h_j,\alpha)$ should price the emitted-key iteration outside heap comparisons and child-mask copies: key construction, fixed state bookkeeping, the directory lookup, and obtaining the current key score or bound. This includes cache effects and collisions, rather than assume one CPU operation.

\begin{align}
T_{C,\rm candidates}={}&T_{\rm keyprep}
 +\sum_{j\in\mathcal P}\Big[
 E_jc_{\rm enum}+C_j^{\rm copy}c_{\rm copy}
 +H_jc_{\rm lookup}(h_j,\alpha)+F_jc_{\rm post}\Big]\notag\\
 &+\sum_j F_jG_{C,j}c_{\rm lut,gather}+H_s(F,R)c_{\rm heap}+T_{\rm rank}(R),
\label{eq:keys-time}\\
T_C={}&\Tcommon+T_{C,\rm candidates}+\Tfinish(R').
\label{eq:keys-total}
\end{align}
$T_{\rm keyprep}$ includes sorting lookup weights, initializing streams and constructing
query keys. For $P$ different free-coordinate sets, a straightforward upper count is
$O(\sum_jh_j\log(h_j+1))$. A global merge adds its own $P$-stream selection work, typically
$O((\sum_jH_j)\log(P+1))$; include it when that variant is used.

The hash directory can be built before the query, but query-specific weighted order depends
on $q$. Precomputing a hash table does not precompute this order for every float query.

\subsection{Query RAM while enumerating}
Let $b_{\rm enum,j}$ include a key mask, penalty and queue metadata. For simultaneously active
streams, the candidate-search scratch is
\begin{equation}
Q_{C,\rm search}=\Qbase+
 \max_t\left[\sum_j\gamma_C Q_j(t)b_{\rm enum,j}
             +Q_{\rm stream\ merge}(t)+Q_{\rm posting\ buffer}(t)
             +Q_{\rm gather}(t)\right].
\label{eq:keys-ram}
\end{equation}
A sequential per-cluster implementation can reuse enumeration memory, replacing the sum of
live stream sizes by the maximum live cluster stream. A global full-score stream may need
all selected streams active. These are different execution choices.

Postings can be streamed into a top-$R$ heap; storing all $F$ retrieved IDs is not mandatory.
If they are materialized for batch reranking, add $F b_{\rm ref}$ or the actual bounded batch
size. A sparse dictionary's small stored size does not bound the size of a query's enumeration
queue. With the two-child stream, the queue can grow by at most one per emitted state, but
$H_j$ itself can be enormous.

\begin{resultbox}{What Approach C does and does not lose}
Full-key enumeration in true penalty order can have 100\% binary recall. Its difficulty can
be time or queue memory, especially when most keys are empty. Short-key search is a different
tradeoff: it can collect useful candidates cheaply, but a fixed quota is not a recall guarantee.
It becomes exact if a valid full-score bound proves that no unseen candidate can improve the
result, or if it eventually scores every eligible document.
\end{resultbox}
